\subsection{Determinant reduction for four-point spectra}
\label{sec:fqe-renyi-determinant}

\begin{lemma}[Determinant monotonicity at fixed second and third moments]
\label{lem:fqe-renyi-determinant}
For $1<q<2$, let $\lambda,\lambda'$ be four-point spectra with the same second moment $p$ and third moment $t$.
If $\prod_i\lambda_i<\prod_i\lambda'_i$, then $H_q(\lambda)<H_q(\lambda')$.
Consequently, at fixed $p$ and $t$, the maximum of $H_q$ is attained at a spectrum of one of the two types
\[
(a,b,c,c)\quad\text{or}\quad(a,a,b,c),
\qquad a\ge b\ge c\ge0,
\tag{DM}\label{eq:fqe-renyi-double-root}
\]
including their degenerations.
\end{lemma}
\begin{proof}[Proof idea]
Write $e_k$ for the elementary symmetric polynomials of the eigenvalues and $d=\prod_i\lambda_i$. Newton's identities fix $e_1,e_2,e_3$, so $\prod_i(\lambda_i+s)$ depends on the spectrum only through $d$; a
Stieltjes representation of $x^q$ then shows that $\sum_i\lambda_i^q$ decreases in $d$, and a spectrum of maximal
determinant has a repeated outer pair. The complete proof is in Section~\ref{sec:sup-determinant} of the supplement~\cite{Companion}.
\end{proof}


\begin{corollary}[Six-variable classical reduction]
\label{cor:fqe-renyi-six-variable}
Let three four-component probability vectors be constrained by any
conditions involving only their total second and third moments $P,T$
and the sum of their squared second moments
$U=\sum_jp_j^2$. For every $1<q<2$, the supremum of their mean
R\'enyi entropy is unchanged if each cut is required to have one
of the spectral forms \eqref{eq:fqe-renyi-double-root}.
Up to permutation of cuts, only four patterns need be considered:
$k=0,1,2,3$ cuts with two largest equal, the others with two
smallest equal. Each cut has two free real spectral parameters.
In particular, this reduction is exact for the classical relaxation
$P\ge1$, $F=17P-12T-12\ge0$, and
$F^2+\tfrac92(P-1)\ge6(U-P^2/3)$.
Theorem~\ref{thm:fqe-renyi-classical} bounds this relaxation at the two anchor orders; with
Corollary~\ref{cor:fqe-renyi-order-transfer} and order monotonicity this gives the envelope on it for $q>1$.
\end{corollary}
\begin{proof}
At fixed $p_j,t_j$ for each cut, replacing that spectrum by one of
maximum determinant preserves $P,T,U$ and thus all the stated
constraints. Lemma~\ref{lem:fqe-renyi-determinant} shows that this
replacement cannot decrease any cut entropy. The two forms in (DM)
have two free coordinates after the trace constraint. The cut
permutation symmetry reduces the eight form choices to four.
\end{proof}
